
%% bare_jrnl.tex
%% V1.3
%% 2007/01/11
%% by Michael Shell
%% see http://www.michaelshell.org/
%% for current contact information.
%%
%% This is a skeleton file demonstrating the use of IEEEtran.cls
%% (requires IEEEtran.cls version 1.7 or later) with an IEEE journal paper.
%%
%% Support sites:
%% http://www.michaelshell.org/tex/ieeetran/
%% http://www.ctan.org/tex-archive/macros/latex/contrib/IEEEtran/
%% and
%% http://www.ieee.org/



% *** Authors should verify (and, if needed, correct) their LaTeX system  ***
% *** with the testflow diagnostic prior to trusting their LaTeX platform ***
% *** with production work. IEEE's font choices can trigger bugs that do  ***
% *** not appear when using other class files.                            ***
% The testflow support page is at:
% http://www.michaelshell.org/tex/testflow/


%%*************************************************************************
%% Legal Notice:
%% This code is offered as-is without any warranty either expressed or
%% implied; without even the implied warranty of MERCHANTABILITY or
%% FITNESS FOR A PARTICULAR PURPOSE! 
%% User assumes all risk.
%% In no event shall IEEE or any contributor to this code be liable for
%% any damages or losses, including, but not limited to, incidental,
%% consequential, or any other damages, resulting from the use or misuse
%% of any information contained here.
%%
%% All comments are the opinions of their respective authors and are not
%% necessarily endorsed by the IEEE.
%%
%% This work is distributed under the LaTeX Project Public License (LPPL)
%% ( http://www.latex-project.org/ ) version 1.3, and may be freely used,
%% distributed and modified. A copy of the LPPL, version 1.3, is included
%% in the base LaTeX documentation of all distributions of LaTeX released
%% 2003/12/01 or later.
%% Retain all contribution notices and credits.
%% ** Modified files should be clearly indicated as such, including  **
%% ** renaming them and changing author support contact information. **
%%
%% File list of work: IEEEtran.cls, IEEEtran_HOWTO.pdf, bare_adv.tex,
%%                    bare_conf.tex, bare_jrnl.tex, bare_jrnl_compsoc.tex
%%*************************************************************************

% Note that the a4paper option is mainly intended so that authors in
% countries using A4 can easily print to A4 and see how their papers will
% look in print - the typesetting of the document will not typically be
% affected with changes in paper size (but the bottom and side margins will).
% Use the testflow package mentioned above to verify correct handling of
% both paper sizes by the user's LaTeX system.
%
% Also note that the "draftcls" or "draftclsnofoot", not "draft", option
% should be used if it is desired that the figures are to be displayed in
% draft mode.
%
\documentclass[journal]{IEEEtran}
\usepackage{graphicx}
\usepackage{lipsum}

% Some very useful LaTeX packages include:
% (uncomment the ones you want to load)


% *** MISC UTILITY PACKAGES ***
%
%\usepackage{ifpdf}
% Heiko Oberdiek's ifpdf.sty is very useful if you need conditional
% compilation based on whether the output is pdf or dvi.
% usage:
% \ifpdf
%   % pdf code
% \else
%   % dvi code
% \fi
% The latest version of ifpdf.sty can be obtained from:
% http://www.ctan.org/tex-archive/macros/latex/contrib/oberdiek/
% Also, note that IEEEtran.cls V1.7 and later provides a builtin
% \ifCLASSINFOpdf conditional that works the same way.
% When switching from latex to pdflatex and vice-versa, the compiler may
% have to be run twice to clear warning/error messages.






% *** CITATION PACKAGES ***
%
%\usepackage{cite}
% cite.sty was written by Donald Arseneau
% V1.6 and later of IEEEtran pre-defines the format of the cite.sty package
% \cite{} output to follow that of IEEE. Loading the cite package will
% result in citation numbers being automatically sorted and properly
% "compressed/ranged". e.g., [1], [9], [2], [7], [5], [6] without using
% cite.sty will become [1], [2], [5]--[7], [9] using cite.sty. cite.sty's
% \cite will automatically add leading space, if needed. Use cite.sty's
% noadjust option (cite.sty V3.8 and later) if you want to turn this off.
% cite.sty is already installed on most LaTeX systems. Be sure and use
% version 4.0 (2003-05-27) and later if using hyperref.sty. cite.sty does
% not currently provide for hyperlinked citations.
% The latest version can be obtained at:
% http://www.ctan.org/tex-archive/macros/latex/contrib/cite/
% The documentation is contained in the cite.sty file itself.






% *** GRAPHICS RELATED PACKAGES ***
%
\ifCLASSINFOpdf
  % \usepackage[pdftex]{graphicx}
  % declare the path(s) where your graphic files are
  % \graphicspath{{../pdf/}{../jpeg/}}
  % and their extensions so you won't have to specify these with
  % every instance of \includegraphics
  % \DeclareGraphicsExtensions{.pdf,.jpeg,.png}
\else
  % or other class option (dvipsone, dvipdf, if not using dvips). graphicx
  % will default to the driver specified in the system graphics.cfg if no
  % driver is specified.
  % \usepackage[dvips]{graphicx}
  % declare the path(s) where your graphic files are
  % \graphicspath{{../eps/}}
  % and their extensions so you won't have to specify these with
  % every instance of \includegraphics
  % \DeclareGraphicsExtensions{.eps}
\fi
% graphicx was written by David Carlisle and Sebastian Rahtz. It is
% required if you want graphics, photos, etc. graphicx.sty is already
% installed on most LaTeX systems. The latest version and documentation can
% be obtained at: 
% http://www.ctan.org/tex-archive/macros/latex/required/graphics/
% Another good source of documentation is "Using Imported Graphics in
% LaTeX2e" by Keith Reckdahl which can be found as epslatex.ps or
% epslatex.pdf at: http://www.ctan.org/tex-archive/info/
%
% latex, and pdflatex in dvi mode, support graphics in encapsulated
% postscript (.eps) format. pdflatex in pdf mode supports graphics
% in .pdf, .jpeg, .png and .mps (metapost) formats. Users should ensure
% that all non-photo figures use a vector format (.eps, .pdf, .mps) and
% not a bitmapped formats (.jpeg, .png). IEEE frowns on bitmapped formats
% which can result in "jaggedy"/blurry rendering of lines and letters as
% well as large increases in file sizes.
%
% You can find documentation about the pdfTeX application at:
% http://www.tug.org/applications/pdftex





% *** MATH PACKAGES ***
%
%\usepackage[cmex10]{amsmath}
% A popular package from the American Mathematical Society that provides
% many useful and powerful commands for dealing with mathematics. If using
% it, be sure to load this package with the cmex10 option to ensure that
% only type 1 fonts will utilized at all point sizes. Without this option,
% it is possible that some math symbols, particularly those within
% footnotes, will be rendered in bitmap form which will result in a
% document that can not be IEEE Xplore compliant!
%
% Also, note that the amsmath package sets \interdisplaylinepenalty to 10000
% thus preventing page breaks from occurring within multiline equations. Use:
%\interdisplaylinepenalty=2500
% after loading amsmath to restore such page breaks as IEEEtran.cls normally
% does. amsmath.sty is already installed on most LaTeX systems. The latest
% version and documentation can be obtained at:
% http://www.ctan.org/tex-archive/macros/latex/required/amslatex/math/





% *** SPECIALIZED LIST PACKAGES ***
%
%\usepackage{algorithmic}
% algorithmic.sty was written by Peter Williams and Rogerio Brito.
% This package provides an algorithmic environment fo describing algorithms.
% You can use the algorithmic environment in-text or within a figure
% environment to provide for a floating algorithm. Do NOT use the algorithm
% floating environment provided by algorithm.sty (by the same authors) or
% algorithm2e.sty (by Christophe Fiorio) as IEEE does not use dedicated
% algorithm float types and packages that provide these will not provide
% correct IEEE style captions. The latest version and documentation of
% algorithmic.sty can be obtained at:
% http://www.ctan.org/tex-archive/macros/latex/contrib/algorithms/
% There is also a support site at:
% http://algorithms.berlios.de/index.html
% Also of interest may be the (relatively newer and more customizable)
% algorithmicx.sty package by Szasz Janos:
% http://www.ctan.org/tex-archive/macros/latex/contrib/algorithmicx/




% *** ALIGNMENT PACKAGES ***
%
%\usepackage{array}
% Frank Mittelbach's and David Carlisle's array.sty patches and improves
% the standard LaTeX2e array and tabular environments to provide better
% appearance and additional user controls. As the default LaTeX2e table
% generation code is lacking to the point of almost being broken with
% respect to the quality of the end results, all users are strongly
% advised to use an enhanced (at the very least that provided by array.sty)
% set of table tools. array.sty is already installed on most systems. The
% latest version and documentation can be obtained at:
% http://www.ctan.org/tex-archive/macros/latex/required/tools/


%\usepackage{mdwmath}
%\usepackage{mdwtab}
% Also highly recommended is Mark Wooding's extremely powerful MDW tools,
% especially mdwmath.sty and mdwtab.sty which are used to format equations
% and tables, respectively. The MDWtools set is already installed on most
% LaTeX systems. The lastest version and documentation is available at:
% http://www.ctan.org/tex-archive/macros/latex/contrib/mdwtools/


% IEEEtran contains the IEEEeqnarray family of commands that can be used to
% generate multiline equations as well as matrices, tables, etc., of high
% quality.


%\usepackage{eqparbox}
% Also of notable interest is Scott Pakin's eqparbox package for creating
% (automatically sized) equal width boxes - aka "natural width parboxes".
% Available at:
% http://www.ctan.org/tex-archive/macros/latex/contrib/eqparbox/





% *** SUBFIGURE PACKAGES ***
%\usepackage[tight,footnotesize]{subfigure}
% subfigure.sty was written by Steven Douglas Cochran. This package makes it
% easy to put subfigures in your figures. e.g., "Figure 1a and 1b". For IEEE
% work, it is a good idea to load it with the tight package option to reduce
% the amount of white space around the subfigures. subfigure.sty is already
% installed on most LaTeX systems. The latest version and documentation can
% be obtained at:
% http://www.ctan.org/tex-archive/obsolete/macros/latex/contrib/subfigure/
% subfigure.sty has been superceeded by subfig.sty.



%\usepackage[caption=false]{caption}
%\usepackage[font=footnotesize]{subfig}
% subfig.sty, also written by Steven Douglas Cochran, is the modern
% replacement for subfigure.sty. However, subfig.sty requires and
% automatically loads Axel Sommerfeldt's caption.sty which will override
% IEEEtran.cls handling of captions and this will result in nonIEEE style
% figure/table captions. To prevent this problem, be sure and preload
% caption.sty with its "caption=false" package option. This is will preserve
% IEEEtran.cls handing of captions. Version 1.3 (2005/06/28) and later 
% (recommended due to many improvements over 1.2) of subfig.sty supports
% the caption=false option directly:
%\usepackage[caption=false,font=footnotesize]{subfig}
%
% The latest version and documentation can be obtained at:
% http://www.ctan.org/tex-archive/macros/latex/contrib/subfig/
% The latest version and documentation of caption.sty can be obtained at:
% http://www.ctan.org/tex-archive/macros/latex/contrib/caption/




% *** FLOAT PACKAGES ***
%
%\usepackage{fixltx2e}
% fixltx2e, the successor to the earlier fix2col.sty, was written by
% Frank Mittelbach and David Carlisle. This package corrects a few problems
% in the LaTeX2e kernel, the most notable of which is that in current
% LaTeX2e releases, the ordering of single and double column floats is not
% guaranteed to be preserved. Thus, an unpatched LaTeX2e can allow a
% single column figure to be placed prior to an earlier double column
% figure. The latest version and documentation can be found at:
% http://www.ctan.org/tex-archive/macros/latex/base/



%\usepackage{stfloats}
% stfloats.sty was written by Sigitas Tolusis. This package gives LaTeX2e
% the ability to do double column floats at the bottom of the page as well
% as the top. (e.g., "\begin{figure*}[!b]" is not normally possible in
% LaTeX2e). It also provides a command:
%\fnbelowfloat
% to enable the placement of footnotes below bottom floats (the standard
% LaTeX2e kernel puts them above bottom floats). This is an invasive package
% which rewrites many portions of the LaTeX2e float routines. It may not work
% with other packages that modify the LaTeX2e float routines. The latest
% version and documentation can be obtained at:
% http://www.ctan.org/tex-archive/macros/latex/contrib/sttools/
% Documentation is contained in the stfloats.sty comments as well as in the
% presfull.pdf file. Do not use the stfloats baselinefloat ability as IEEE
% does not allow \baselineskip to stretch. Authors submitting work to the
% IEEE should note that IEEE rarely uses double column equations and
% that authors should try to avoid such use. Do not be tempted to use the
% cuted.sty or midfloat.sty packages (also by Sigitas Tolusis) as IEEE does
% not format its papers in such ways.


%\ifCLASSOPTIONcaptionsoff
%  \usepackage[nomarkers]{endfloat}
% \let\MYoriglatexcaption\caption
% \renewcommand{\caption}[2][\relax]{\MYoriglatexcaption[#2]{#2}}
%\fi
% endfloat.sty was written by James Darrell McCauley and Jeff Goldberg.
% This package may be useful when used in conjunction with IEEEtran.cls'
% captionsoff option. Some IEEE journals/societies require that submissions
% have lists of figures/tables at the end of the paper and that
% figures/tables without any captions are placed on a page by themselves at
% the end of the document. If needed, the draftcls IEEEtran class option or
% \CLASSINPUTbaselinestretch interface can be used to increase the line
% spacing as well. Be sure and use the nomarkers option of endfloat to
% prevent endfloat from "marking" where the figures would have been placed
% in the text. The two hack lines of code above are a slight modification of
% that suggested by in the endfloat docs (section 8.3.1) to ensure that
% the full captions always appear in the list of figures/tables - even if
% the user used the short optional argument of \caption[]{}.
% IEEE papers do not typically make use of \caption[]'s optional argument,
% so this should not be an issue. A similar trick can be used to disable
% captions of packages such as subfig.sty that lack options to turn off
% the subcaptions:
% For subfig.sty:
% \let\MYorigsubfloat\subfloat
% \renewcommand{\subfloat}[2][\relax]{\MYorigsubfloat[]{#2}}
% For subfigure.sty:
% \let\MYorigsubfigure\subfigure
% \renewcommand{\subfigure}[2][\relax]{\MYorigsubfigure[]{#2}}
% However, the above trick will not work if both optional arguments of
% the \subfloat/subfig command are used. Furthermore, there needs to be a
% description of each subfigure *somewhere* and endfloat does not add
% subfigure captions to its list of figures. Thus, the best approach is to
% avoid the use of subfigure captions (many IEEE journals avoid them anyway)
% and instead reference/explain all the subfigures within the main caption.
% The latest version of endfloat.sty and its documentation can obtained at:
% http://www.ctan.org/tex-archive/macros/latex/contrib/endfloat/
%
% The IEEEtran \ifCLASSOPTIONcaptionsoff conditional can also be used
% later in the document, say, to conditionally put the References on a 
% page by themselves.





% *** PDF, URL AND HYPERLINK PACKAGES ***
%
%\usepackage{url}
% url.sty was written by Donald Arseneau. It provides better support for
% handling and breaking URLs. url.sty is already installed on most LaTeX
% systems. The latest version can be obtained at:
% http://www.ctan.org/tex-archive/macros/latex/contrib/misc/
% Read the url.sty source comments for usage information. Basically,
% \url{my_url_here}.





% *** Do not adjust lengths that control margins, column widths, etc. ***
% *** Do not use packages that alter fonts (such as pslatex).         ***
% There should be no need to do such things with IEEEtran.cls V1.6 and later.
% (Unless specifically asked to do so by the journal or conference you plan
% to submit to, of course. )


% correct bad hyphenation here
\hyphenation{op-tical net-works semi-conduc-tor}


\begin{document}
%
% paper title
% can use linebreaks \\ within to get better formatting as desired
%\title{Post Quantum Cryptography}
\title{The Impact of Quantum Computing on Cryptography}
%
%
% author names and IEEE memberships
% note positions of commas and nonbreaking spaces ( ~ ) LaTeX will not break
% a structure at a ~ so this keeps an author's name from being broken across
% two lines.
% use \thanks{} to gain access to the first footnote area
% a separate \thanks must be used for each paragraph as LaTeX2e's \thanks
% was not built to handle multiple paragraphs
%

\author{Tara O'Kelly,~\IEEEmembership{(Hons) Software Development,~GMIT}% <-this % stops a space
}

% note the % following the last \IEEEmembership and also \thanks - 
% these prevent an unwanted space from occurring between the last author name
% and the end of the author line. i.e., if you had this:
% 
% \author{....lastname \thanks{...} \thanks{...} }
%                     ^------------^------------^----Do not want these spaces!
%
% a space would be appended to the last name and could cause every name on that
% line to be shifted left slightly. This is one of those "LaTeX things". For
% instance, "\textbf{A} \textbf{B}" will typeset as "A B" not "AB". To get
% "AB" then you have to do: "\textbf{A}\textbf{B}"
% \thanks is no different in this regard, so shield the last } of each \thanks
% that ends a line with a % and do not let a space in before the next \thanks.
% Spaces after \IEEEmembership other than the last one are OK (and needed) as
% you are supposed to have spaces between the names. For what it is worth,
% this is a minor point as most people would not even notice if the said evil
% space somehow managed to creep in.



% The paper headers
\markboth{The Impact of Quantum Computing on Cryptography, October~2017}%
{The Impact of Quantum Computing on Cryptography, October~2017}
% The only time the second header will appear is for the odd numbered pages
% after the title page when using the twoside option.
% 
% *** Note that you probably will NOT want to include the author's ***
% *** name in the headers of peer review papers.                   ***
% You can use \ifCLASSOPTIONpeerreview for conditional compilation here if
% you desire.




% If you want to put a publisher's ID mark on the page you can do it like
% this:
%\IEEEpubid{0000--0000/00\$00.00~\copyright~2007 IEEE}
% Remember, if you use this you must call \IEEEpubidadjcol in the second
% column for its text to clear the IEEEpubid mark.



% use for special paper notices
%\IEEEspecialpapernotice{(Invited Paper)}



% https://tex.stackexchange.com/questions/140579/ieee-tran-one-column-two-column-elements-in-the-same-page-possible

\twocolumn[
\begin{@twocolumnfalse}
\maketitle
\end{@twocolumnfalse}
\begin{abstract}
%The security implications of quantum computers has seen to the investigation of quantum resistant cryptography. 
With the arrival of quantum computing, one must acknowledge the security implications. The integrity of many modular cryptographic systems, namely discrete logarithm and factoring based systems, are at risk. Examining the pre-quantum state of cryptography, we will discuss the predicted affect of quantum on cryptography.
%and potential collateral ramifications. 
%Subsequently, we will delve into viable post-quantum solutions. 
Although there are quantum computers that exist today, all are far from capable of performing operations complex enough to break cyptographic algorithms widely used today. 
This does not mean an impact has not yet been felt. The threat of quantum computing has already lead to the creation a new cryptographic field - post-quantum crytography.
The attempts of this field are in contention with the emerging technology, pursuing to implement a \textbf{\textit{feasible}}, \textbf{\textit{flexible}} and \textbf{\textit{computationally efficient}} solution before a pertinent quantum computer can be built.

%With the arrival of quantum computing, one must acknowledge the security implications of quantum computers. The integrity of many modular cryptographic systems, namely discrete logarithm and factoring based systems \cite{Sec_Risk}, are at risk. Examining the pre-quantum state of cryptography and the predicted affect of quantum on cryptography, we will discuss the need for post-quantum computing. Subsequently, we will delve into the viable solutions. Although there are quantum computers that exist today, all are far from capable of performing operations complex enough to break cyptographic algorithms widely used today. The attempts of post-quantum cryptography are in contention with the emerging technology, pursuing to implement a \textbf{\textit{feasible}}, \textbf{\textit{flexible}} and \textbf{\textit{computationally efficient}} solution before a pertinent quantum computer can be built.
\end{abstract}
%\begin{IEEEkeywords}
%IEEEtran, journal, \LaTeX, paper, template.
%\end{IEEEkeywords}
\bigskip]
% IEEEtran.cls defaults to using nonbold math in the Abstract.
% This preserves the distinction between vectors and scalars. However,
% if the journal you are submitting to favors bold math in the abstract,
% then you can use LaTeX's standard command \boldmath at the very start
% of the abstract to achieve this. Many IEEE journals frown on math
% in the abstract anyway.

% Note that keywords are not normally used for peerreview papers.






% For peer review papers, you can put extra information on the cover
% page as needed:
% \ifCLASSOPTIONpeerreview
% \begin{center} \bfseries EDICS Category: 3-BBND \end{center}
% \fi
%
% For peerreview papers, this IEEEtran command inserts a page break and
% creates the second title. It will be ignored for other modes.
\IEEEpeerreviewmaketitle


\section{Introduction}
%The continuous growth in the industry of technology reaps many upsides. 
Quantum computing algorithms have been found to solve certain complex problems abundantly faster than that of classical algorithms, benefiting various areas like machine learning \cite{Sec_Risk} and developing new drugs \cite{Q_Comp_Real}. 
%However, crossing new boundaries can also create downsides. 
However, these algorithms could also aid the devaluation of cryptographic systems in our digital entrenched world.
%RSA \& ECC v popular. Warn about the following being a basic review to understand the relevant points. With benefits come repercussions.
% Cryptography has been around for thousands of years. Nowadays, is at the heart of worldwide communication network.



The following review is not meant to provide in-depth understanding quantum computing. Nor is it intended to be a thorough explanation of current cryptographic algorithms, but rather a basic understanding of their workings and where quantum algorithms compromise them.
Subsequently, we will delve into viable post-quantum solutions proposed in attempts to conserve secure communications.

%\subsection{The Art of Encryption}
\section{Introduction to Cryptography}
Cryptography is a method of protecting data from people who are not authorized to access it. Encryption, a significant mechanism in cryptography, is achieved by transforming plaintext into ciphertext, with the intent of rendering it meaningless to those who do not have the classified resources. With these resources, the ciphertext is usually transformed it back to it's original state, a process called decryption.

Symmetric-key encryption utilizes the same key to both encrypt and decrypt, Advanced Encryption Standard (AES) being the prevailing form. Symmetric encryption is more accomplished in achieving a faster outcome than that of asymmetric. Problematically, the sender and the receiver must hold the same key. How do they exchange keys digitally without another party eavesdropping? A typical solution would be to use asymmetric encryption to exchange symmetric keys.

Asymmetric-key encryption (a.k.a. public-key cryptography) practices an encryption technique with two differing, mathematically linked keys. The first key being a public key and the corresponding key being a private key  - e.g. Rivest Shamir Adleman (RSA), Elliptic Curve Cryptography (ECC). As implied in it's name, the public key is to be publicized. The private key is to be kept secret. The public key can be used to encrypt data, whereas the homologous private key can be used to decrypt the data. An ideal public-key cryptographic algorithm should serve as a trapdoor function. A trapdoor function is a function that is easy to perform one way, but has a secret that is required to perform the inverse calculation efficiently. The objective is to decrease the probability that the secret could be identified as much as possible.

So far, we have acknowledged AES, RSA and ECC; all can be adopted in pursual of a \textbf{confidential} message. Cryptography is also concerned with proving the \textbf{integrity}, \textbf{authenticity} and \textbf{non-repudiation} of a message e.g. Message Authentication Code (MAC), Digital Signatures. Another method attempting the assurance of integrity could be by hash functions, e.g. Secure Hash Algorithm (SHA).

\subsection{Rivest Shamir Adleman}
 Rivest Shamir Adleman was first publicized in 1978 by Ron Rivest, Adi Shamir and Leonard Adleman. The RSA algorithm is the most popular and  perhaps the best understood public key cryptography system. RSA's security derives from the difficulty of factoring two large integers and the expeditious multiplication to get these large numbers; it is a befitting example of the ``trapdoor" methodology. 
 
 Two prime numbers, \(p\) and \(q\), are generated. The values \(p\) and \(q\) are multiplied together to get the maximum value, \(n\). 
 %A number \(pub\) is selected, such that \(pub\) is not a factor of \((p - 1)\) and \((q - 1)\). Then \(priv\) is generated, such that \((priv * pub) mod (p - 1) * (q - 1) = 1\). 
 Born from these values are the numbers \(pub\) and \(priv\).
 The public key \((pub,n)\) can be released, but the private key \((priv,p,q)\) must be kept confidential.
 %The message \(M\) is converted to an integer \(C\) by multiplying itself to the power of the public key, then employs a wrapping scheme to ensure \(C > 0\) and \(C < n\). The decryption is a similar process, the cipher \(C\) is converted to \(M\) by squaring itself to the power of the private key. A wrapping scheme is again implemented, ensuring \(M > 0\) and \(M < n\). 
 As long as you know the values \(p\) and \(q\), you can compute the corresponding private key from the public key, explaining how factoring relates to breaking RSA. Factoring the maximum number into its component primes allows you to compute someone's private key from the public key and decrypt their private messages.

\subsection{Elliptic Curve Cryptography}

Although Elliptic Curve Cryptography was originally proposed in 1985 by by Neal Koblitz and Victor S. Miller, it was not widely utilized until the 21st century. It is not as widely understood as RSA, with it's complexity to blame. ECC allows the use of smaller keys than RSA to get the same levels of security. 
%Small keys are very beneficial nowadays, considering cryptography is often implemented on low powered devices, e.g. a mobile phone. While multiplying two prime numbers together is easier, when the prime numbers start to get very long, even just the multiplication step can take some time on a low powered device.

ECC is based on the algebraic structure of elliptic curves over finite fields \cite{RSA_ECC}. ECC handles the following domain parameters: \((p,a,b,G,n,h)\). To briefly explain the parameters, \(p\) is the field that the graph is defined over, the variables \(a\) and \(b\) are values that define the curve, \(G\) is known as the generator point (a.k.a. base point), \(n\) is the prime order of G and \(h\) is the cofactor of the curve.

The private key \(d\) is a randomly selected integer in the interval \( [ 1 , n - 1 ]\). Subsequently, the public key \(Q=dG\). The security of ECC is built upon the Elliptic Curve Discrete Logarithm Problem (ECDLP) \cite{RSA_ECC}; ECDLP in ECC applies to the laborious task of locating the discrete logarithm of random elliptic curve element even with a known point. With the given the domain parameters and \(Q\), ECDLP refers to the problem of determining \(d\). In a real world standard application, it would be unfeasible to check all the possibilities of \(d\).
\section{Quantum Computing}
Quantum computers apply quantum mechanics to perform computations. Quantum mechanics is the study of the laws of nature on an atomic and subatomic level. The behaviour at this scaled-down level can be replicated at certain temperatures; temperatures just above absolute zero as reported by Conover \cite{Q_Comp_Real}.

Unlike the binary bits of a classical computer, qubits utilize subatomic particles to represent data. 
This allows qubits the additional possibility of being in more than one state simultaneously.
A phenomenon known as a superposition. Therefore, a qubit can be in the state of either 0, 1 or a combination of both. Once a qubit is read, it is similar to a classical bit in that it can only be one of two states, 0 or 1.

Given a superposition, all states of the qubit can be operated on at the same time. This is called quantum parallelism and induces faster computations. So on an \(n\)-qubit computer, \(2n\) values can be computed \cite{Impact_QC_Cryptog}. The speed of a quantum computer can also be credited to it's entangled state \cite{Impact_QC_Cryptog_1}. 
Qubits can be entangled, a sort of relationship where the state of one qubit depends on the state of another. For instance, this enables us to determine the state of a qubit in an entangled pair by observing the state of the other qubit. 
However, the exploitation of quantum behaviour does not expedite the performance of all tasks \cite{Q_Comp_Real}. Only certain quantum algorithms for specific tasks are capable of surpassing the efficiency of classical algorithms. To date, the algorithms for factoring and searching by Shor and Grover are the most promising \cite{Impact_QC_Cryptog_1}.

\subsection{Shor's Algorithm}
Shor's Algorithm has a severe impact on cryptography, potentially breaking the public key cryptography systems used today \cite{Q_Alg}.
In 1994, Peter Shor established a quantum algorithm to find the prime factors of an integer in polynomial time (\(O(N^c)\) time) as opposed to the exponential time (\(O(c^N)\) time) of a classical computer with the same objective. 
The proficient factorization can be obtained as all of the values of a certain function can be computed by applying quantum parallelism and entanglement \cite{Impact_QC_Cryptog_1}. 
%Wright - "quantum alg has been dev to find the prime factors of an integer in polynomial time rather than exponential time. It uses quantum parallelism and entanglement to simultaneously compute all of the values of a certain function."
Shor has also released an adapted version thus damaging the impregnability of the discrete logarithm problem and ECDLP \cite{Post_Q_Cryptog}. 

Public key encryption and digital signatures such as the favored algorithms RSA, ECC, Digital Signature Algorithm and Elliptic Curve Digital Signature Algorithm are vulnerable in this scenario. With the corruption of public key encryption comes the simultaneous weakening of single key encryption by straining the means to disclose symmetric keys.

\subsection{Grover's Algorithm}
Grover's Algorithm was formulated by Lov Grover in 1996. This algorithm is a strong foundation on which lot of the superior applications of quantum computing can be built \cite{Post_Q_Cryptog}. Harmful affects by the algorithm in question can be seen in symmetric encryption and hash functions, including the ever popular AES and the SHA family. It is important to note that Grover's algorithm only decreases the security of these cryptographic methods; it does not break them.

A classical computer can do a search with the speed of \(O(N)\) assuming we do not know if the search parameter exists in an unordered domain. Grover's algorithm suggests that the same search can be implemented in \(O(\sqrt{N}\)) time using quantum queries. Grover's acceleration from \(O(N)\) to \(O(\sqrt{N}\)) is not quite as calamitous as Shor's. However, it indicates another significant quantum approach to cryptanalysis \cite{Quantum_Cryptanal}, the study of code breaking. 

\subsection{Security Implications}
The exertion of the quantum algorithms above on a capable quantum computer has the power to revolutionize cryptanalysis. Fortunately, no quantum computer is qualified to do so yet. As Bacon and Dam stated, ``We can communicate securely, today, given that we cannot build a large scale quantum computer tomorrow" \cite{Q_Alg}. This extensive advancement in cryptanalysis will presumably open the door to a copious amount of privacy issues. 

Under the assumption of a relevant quantum computer, an unauthorized third party could eavesdrop on transactions and even manipulate said transactions. Additionally, this third party could easily impersonate others - even trusted sites, with the debilitation of digital certificates \cite{Sec_Risk}. The cryptographic principles confidence, integrity, authentication and non-repudiation would all be put in jeopardy. Many types of digital transactions would be called into question; personal banking, online purchases and software downloads/updates to name a few. Not to mention the immense volume of private information accumulated in our cloud based world, such as the vast amount of personal information submitted to social networks and held in medical records, would no longer be secure. 

The variety of reasons to motivate malicious hackers is extensive. For example, a company may be seeking to gain an unfair advantage over competitors or a government may wish to monitor transactions to identify and combat threats. The numerous different possibilities of hackers, motivations and victims each spark worries over any conceivable collateral ramifications. A government's acquisition of a quantum cryptanalysis system can be seen as a particularly compelling concern \cite{Sec_Risk}. The exploitation of citizen privacy may occur. The NSA's leaked documents revealed PRISM, a strict surveillance program, that can be referenced in this case \cite{Sec_Risk}. PRISM achieved extreme monitoring with the cooperation of service providers. A program such as PRISM can lead to a harmful surveillance society, possibly resulting in citizen unrest and mistrust of their governing authority. In a drastic instance, global tensions could be escalated by increasing a governments ability to spy on another. Whether it be by discovering another administration's private information, or the detection of another administration observing their state or their civilian's private information. The very same leaked documents also exposed a plan to construct a cryptologically useful quantum computer.
Of course it is not only powerful entities, like a sizable company or a governing body, with the means to develop or obtain a quantum computer that are applicable to be considered a threat. One does not have to own a quantum computer, but could rely on a service provided by another party that possesses a quantum computer. 

\subsection{Post-Quantum Cryptography}
The security implications of quantum computers has seen to the investigation of quantum resistant cryptography. The less worrisome weakening of symmetric key cryptography can be addressed by increasing the security level. If not doing so already, Bernstein and Lange recommend simply switching to AES 256-bit keys \cite{Post_Q_Cryptog}. 

Tackling asymmetric cryptography is another matter.
Given that no relevant computer has been made yet, hope is not lost. 
Cryptographers worldwide have been attempting to find a solution to this imminent threat. 
Both classical and quantum solutions have been explored. Many solutions have been proposed, but only a handful of these have persevered through the extensive scrutiny.
It is difficult to actuate the strength of these proposals without the meticulous means to test them, e.g. a quantum computer at sufficient potential, time. 

%As synopsized by Bernstein and Lange \cite{Post_Q_Cryptog}, code-based encryption offers high confidence, fast encryption and short outputs. Unfortunately, the public keys are terribly large.
%The speed of lattice-based encryption/signatures are also favourable. Both short keys and outputs are produced. However, lattice-based cryptography requires more security analysis.
%Similarly, multivariate-quadratic-equation signatures grant short outputs, but also needs more security analysis.
As synopsized by Bernstein and Lange, code-based encryption, lattice-based encryption/signatures, multivariate-quadratic-equation signatures and hash-based signatures are practical classical proposals \cite{Post_Q_Cryptog}. However, none of these classical proposals are without shortcomings. For example, the likes of code-based encryption produces terribly large public keys \cite{Post_Q_Cryptog}, hash-based signatures are burdened with large signatures \cite{Sec_Risk} and as some of them are quite new, they are in need of additional time to further endure lengthy cryptanalysis \cite{Sec_Risk}.


%Hash-based signatures
Ironically, the same technology expected to suppress current security techniques also harbors great promise in bringing forth strong replacements \cite{Q_Cryptog}. 
%Most well studied is the Quantum Key Distribution Protocol (QKDP). 
Quantum Key Distribution (QKD) protocols make use of photons (particles of light) as opposed to computational complexity to distribute symmetric keys \cite{Impact_QC_Cryptog_1}. QKD is already commercially available, but only over limited distances, e.g. ID Quantique \cite{IDQ}. The first QKD protocol was published in 1984 by Bennett and Brassard, respectively earning the name BB84. Thereafter came Ekert's similarly named E91 and Bennett's solo B92 protocol \cite{Q_Cryptog}.


% needed in second column of first page if using \IEEEpubid
%\IEEEpubidadjcol

% An example of a floating figure using the graphicx package.
% Note that \label must occur AFTER (or within) \caption.
% For figures, \caption should occur after the \includegraphics.
% Note that IEEEtran v1.7 and later has special internal code that
% is designed to preserve the operation of \label within \caption
% even when the captionsoff option is in effect. However, because
% of issues like this, it may be the safest practice to put all your
% \label just after \caption rather than within \caption{}.
%
% Reminder: the "draftcls" or "draftclsnofoot", not "draft", class
% option should be used if it is desired that the figures are to be
% displayed while in draft mode.
%
%\begin{figure}[!t]
%\centering
%\includegraphics[width=2.5in]{myfigure}
% where an .eps filename suffix will be assumed under latex, 
% and a .pdf suffix will be assumed for pdflatex; or what has been declared
% via \DeclareGraphicsExtensions.
%\caption{Simulation Results}
%\label{fig_sim}
%\end{figure}

% Note that IEEE typically puts floats only at the top, even when this
% results in a large percentage of a column being occupied by floats.


% An example of a double column floating figure using two subfigures.
% (The subfig.sty package must be loaded for this to work.)
% The subfigure \label commands are set within each subfloat command, the
% \label for the overall figure must come after \caption.
% \hfil must be used as a separator to get equal spacing.
% The subfigure.sty package works much the same way, except \subfigure is
% used instead of \subfloat.
%
%\begin{figure*}[!t]
%\centerline{\subfloat[Case I]\includegraphics[width=2.5in]{subfigcase1}%
%\label{fig_first_case}}
%\hfil
%\subfloat[Case II]{\includegraphics[width=2.5in]{subfigcase2}%
%\label{fig_second_case}}}
%\caption{Simulation results}
%\label{fig_sim}
%\end{figure*}
%
% Note that often IEEE papers with subfigures do not employ subfigure
% captions (using the optional argument to \subfloat), but instead will
% reference/describe all of them (a), (b), etc., within the main caption.


% An example of a floating table. Note that, for IEEE style tables, the 
% \caption command should come BEFORE the table. Table text will default to
% \footnotesize as IEEE normally uses this smaller font for tables.
% The \label must come after \caption as always.
%
%\begin{table}[!t]
%% increase table row spacing, adjust to taste
%\renewcommand{\arraystretch}{1.3}
% if using array.sty, it might be a good idea to tweak the value of
% \extrarowheight as needed to properly center the text within the cells
%\caption{An Example of a Table}
%\label{table_example}
%\centering
%% Some packages, such as MDW tools, offer better commands for making tables
%% than the plain LaTeX2e tabular which is used here.
%\begin{tabular}{|c||c|}
%\hline
%One & Two\\
%\hline
%Three & Four\\
%\hline
%\end{tabular}
%\end{table}


% Note that IEEE does not put floats in the very first column - or typically
% anywhere on the first page for that matter. Also, in-text middle ("here")
% positioning is not used. Most IEEE journals use top floats exclusively.
% Note that, LaTeX2e, unlike IEEE journals, places footnotes above bottom
% floats. This can be corrected via the \fnbelowfloat command of the
% stfloats package.



\section{Conclusion}
%It is quite apparent that quantum computers has already left it's mark on cryptography. And if not challenged, could turn that mark into a deep wound.
According to Kudelski there is no need to worry about quantum computers breaking current cryptographic systems for years to come; cryptographers have time to devise an optimal solution \cite{Impact_QC_Cryptog}. However, one must also consider the time costs of integration and standardization. It could take companies years to fully update their systems. Classical post-quantum solutions are complex, therefore costly to implement \cite{Sec_Risk}. Due to the uncertainty of quantum computers, entities are understandably hesitant to hastily invest in a new security system. And for quantum cryptography to become commercially viable, the quantum internet would have to be more efficient and expansive. This depends on the advances in physics and technology \cite{Q_Cryptog}. But with a quantum satellite already orbiting the earth, the eventual future does look prosperous.

So can our security systems be quantum safe on time? Post-quantum cryptography is fighting to meet an unknown deadline, so there is no way to be sure. In fact, this review has been written under the assumption that the state of quantum computing is as publicly known. There is always the possibility that a capable quantum computer may be closer than we think. Not without regarding the extreme unlikeliness of the following, one may already exist. 




% if have a single appendix:
%\appendix[Proof of the Zonklar Equations]
% or
%\appendix  % for no appendix heading
% do not use \section anymore after \appendix, only \section*
% is possibly needed

% use appendices with more than one appendix
% then use \section to start each appendix
% you must declare a \section before using any
% \subsection or using \label (\appendices by itself
% starts a section numbered zero.)
%


%\appendices
%\section{Proof of the First Zonklar Equation}
%Some text for the appendix. This is obvious.

% use section* for acknowledgement
%\section*{Acknowledgment}


%The authors would like to thank...


% Can use something like this to put references on a page
% by themselves when using endfloat and the captionsoff option.
\ifCLASSOPTIONcaptionsoff
  \newpage
\fi



% trigger a \newpage just before the given reference
% number - used to balance the columns on the last page
% adjust value as needed - may need to be readjusted if
% the document is modified later
%\IEEEtriggeratref{8}
% The "triggered" command can be changed if desired:
%\IEEEtriggercmd{\enlargethispage{-5in}}

% references section

% can use a bibliography generated by BibTeX as a .bbl file
% BibTeX documentation can be easily obtained at:
% http://www.ctan.org/tex-archive/biblio/bibtex/contrib/doc/
% The IEEEtran BibTeX style support page is at:
% http://www.michaelshell.org/tex/ieeetran/bibtex/
%\bibliographystyle{IEEEtran}
% argument is your BibTeX string definitions and bibliography database(s)
%\bibliography{IEEEabrv,../bib/paper}
%
% <OR> manually copy in the resultant .bbl file
% set second argument of \begin to the number of references
% (used to reserve space for the reference number labels box)

\bibliography{references} 
\bibliographystyle{IEEEtran}

% biography section
% 
% If you have an EPS/PDF photo (graphicx package needed) extra braces are
% needed around the contents of the optional argument to biography to prevent
% the LaTeX parser from getting confused when it sees the complicated
% \includegraphics command within an optional argument. (You could create
% your own custom macro containing the \includegraphics command to make things
% simpler here.)
%\begin{biography}[{\includegraphics[width=1in,height=1.25in,clip,keepaspectratio]{mshell}}]{Michael Shell}
% or if you just want to reserve a space for a photo:

%\begin{IEEEbiography}[{\includegraphics[width=1in,height=1.25in,%clip,keepaspectratio]{picture}}]{John Doe}
%\blindtext
%\end{IEEEbiography}

% You can push biographies down or up by placing
% a \vfill before or after them. The appropriate
% use of \vfill depends on what kind of text is
% on the last page and whether or not the columns
% are being equalized.

%\vfill

% Can be used to pull up biographies so that the bottom of the last one
% is flush with the other column.
%\enlargethispage{-5in}



% that's all folks
\end{document}


